The \textsc{bisect-apportion} crate establishes that a redistricting
platform can implement federal apportionment law exactly, verifiably,
and with full test coverage.
The SHA-256 verification protocol is the key contribution:
it creates a cryptographic link between the Census Bureau's official
publication and the redistricting platform's apportionment step,
providing the auditability that litigation and commission contexts require.

Three engineering decisions underpin the crate's correctness:
\begin{enumerate}
  \item \textbf{Exact integer arithmetic}: The \texttt{u128} rational
    representation for priority values eliminates floating-point
    rounding errors in the tie-breaking step, which is where naive
    implementations are most likely to produce incorrect results.
  \item \textbf{Unified architecture}: The shared priority-queue core
    for all divisor methods ensures that correctness demonstrated for
    Huntington-Hill (via SHA verification) transfers structurally
    to Webster, Adams, and Jefferson.
  \item \textbf{L0/L1/L2 test pyramid}: The 61-test suite covers
    unit formulas (L0), synthetic examples (L1), and real-data regression
    (L2), providing the multilevel assurance expected for software
    implementing federal law.
\end{enumerate}

The comparative analysis (Section~\ref{sec:comparison}) shows that for
the 2020 Census, Huntington-Hill, Webster, and Hamilton agree on all
50 states, while Adams and Jefferson produce small divergences for
states near the rounding boundary.
This empirical result contextualises the theoretical debates in J.1--J.5:
the choice of apportionment method has real but limited consequences
for 2020 Census data, concentrated in a small number of states with
populations near the rounding threshold.

For practitioners, the key takeaway is that \textsc{bisect-apportion}
provides a legally defensible, computationally exact, and independently
auditable implementation of 2~U.S.C.\ \S2a.
It should be the first step in any redistricting pipeline:
run \texttt{bisect apportion --year \{year\} --method huntington-hill
--verify} before beginning any district-drawing computation,
and the SHA verification will confirm that the seat counts are correct.

Future work should extend the verification to state legislative
apportionments (which may use different methods and different
statutory seat counts) and to international PR seat allocations
(using the d'Hondt and Sainte-Lagu\"{e} equivalences documented in J.2 and J.4).
